We have presented a systematic first-principles computational study
of band-gap evolution with temperature in three lead-free tin-halide
perovskites --- \CsSnI{}, \FASnI{}, and \MASnI{} --- using a
validated machine-learning force field for finite-temperature
molecular dynamics and an accelerated SS-DFT framework for
band-gap evaluation across large supercells.
By separately quantifying the contributions of lattice expansion,
polymorphic phase identity, and organic-cation orientational disorder,
we have shown that the net temperature dependence of band gap is
a composite effect that cannot be attributed to any single structural
mechanism.

The principal findings are as follows.

\begin{enumerate}
  \item \textbf{Lattice expansion consistently opens the band gap.}
        Across all studied materials and phases, increasing the cell
        volume raises the band gap, with a sensitivity of
        approximately \SI{0.075}{eV} per \SI{1}{\percent} volume
        expansion.
        This is the most transferable quantitative relationship
        identified in this study.

  \item \textbf{Phase identity independently shifts the band-gap baseline.}
        Different polymorphs of the same composition do not share
        the same band-gap baseline at equilibrium volume.
        The magnitude of the phase-dependent shift is
        composition-dependent: comparatively small for \CsSnI{} and
        more pronounced for the hybrid systems \FASnI{} and \MASnI{}.

  \item \textbf{Organic-cation orientation affects band gap in a
        cation- and phase-specific way.}
        In the tetragonal phase, \MASnI{} is far more sensitive to
        cation rotation than \FASnI{}; in the high-temperature cubic
        phase, this sensitivity reverses.
        The MA/FA contrast demonstrates that A-site cation chemistry
        is not a simple scalar perturbation on band gap but a
        structurally coupled degree of freedom whose influence depends
        critically on the surrounding polymorph.

  \item \textbf{Under combined finite-temperature and strain conditions,}
        the all-inorganic \CsSnI{} retains a regular and predictable
        strain response, while the hybrid systems show complex,
        phase-dependent behaviour that reflects the competing action
        of strain and cation disorder.

  \item \textbf{The established computational workflow} --- comprising
        MLFF preparation, long-time finite-temperature MD, and
        SS-DFT band-gap evaluation on large supercells ---
        successfully reproduces the orthorhombic-to-tetragonal
        transition in the hybrid systems and is directly applicable
        to a broader class of lead-free perovskite materials.
\end{enumerate}

These results provide a mechanistic foundation for rational
band-gap engineering in lead-free tin-halide perovskites.
They identify A-site cation identity and orientational dynamics as
an underexplored control variable in thermal band-gap management,
and they highlight the importance of separately accounting for
phase evolution and cation disorder when interpreting experimental
temperature-dependent measurements.
Future work should extend the cation-disorder analysis to fully
converged finite-temperature ensembles, include spin--orbit coupling
for quantitative comparison with experiment, and extend the integrated
temperature-strain workflow to surfaces, interfaces, and
defect-containing structures, where manufacturing-relevant effects
are most directly accessible.